% =====================================================================
% TEMPLATE — Detailed algorithm description (methodology asset A-algdoc-tpl)
% =====================================================================
% Purpose: the TEMPLATE LAYER of the "document algorithms in detail" convention.
% Any algorithm phase — P5 (MILP formulation), P6 (MILP implementation), P8
% (heuristic design) — must produce a DEDICATED, detailed write-up of its
% algorithm (not merely the one-paragraph P10 report summary) so the technique
% can be reproduced, reviewed and handed off. This skeleton is that write-up.
%
% First instance: docs/<problem>/heuristic.tex (a full Iterated Greedy
% write-up). Reuse this file for the MILP formulation/implementation docs too.
%
% Fill every <PLACEHOLDER>; delete guidance comments. Meant to be \input from a
% main.tex, whose preamble provides at least:
%     \usepackage{amsmath,amssymb}
% (see the working harness in docs/<problem>/ — main.tex + compile.py).
% State in the header comment WHICH phase this doc covers and WHERE it slots in
% the \input order (e.g. after the formulation part, before the report).
%
% Structure (keep the skeleton stable; enrich guidance only):
%   Overview → per-stage/component sections (each with verbatim pseudocode) →
%   Core routine / objective evaluator → Complexity → Parameters (table) →
%   Worked micro-example → Design rationale & limitations.
%
% Principles this doc must honor:
%   * Pseudocode is verbatim and self-contained — a reader reproduces the stage
%     from it alone, without the source code.
%   * Name the CORE reusable routine explicitly (the objective evaluator / cost
%     function most other stages call), and say why it is checker-valid /
%     correct-by-construction where that applies.
%   * State complexity per stage AND overall, in the algorithm's own size
%     symbols (define them once).
%   * Every tunable knob appears in the Parameters table with a DEFAULT.
%   * The micro-example is small enough to trace by hand end to end.
%   * Rationale explains WHY this design (tie back to the requirements / the
%     method-review choice); limitations are honest, not decorative.
% =====================================================================

% =====================================================================
\part{<ALGORITHM NAME> --- detailed description}
% =====================================================================

\section{Overview}
% 4–8 sentences. Name the algorithm and, if it came from a method-review phase,
% the candidate/source it realizes. State the objective it optimizes (the single
% headline metric). List its stages/components as an enumerate, in execution
% order, and point at the file(s) that implement it. End by naming the CORE
% reusable routine (the objective evaluator / cost function) that ties the
% pipeline together and makes this technique comparable to the others.
<OVERVIEW: what the algorithm is, what it optimizes, and its stages in order>

The algorithm has <N> stages, implemented in \texttt{<path/to/impl>}:
\begin{enumerate}
  \item \textbf{<STAGE 1 NAME>} --- <one line: what it decides / produces>.
  \item \textbf{<STAGE 2 NAME>} --- <one line>.
  \item \textbf{<STAGE 3 NAME>} --- <one line>.
  % add/remove stages to match the algorithm
\end{enumerate}

% One sentence naming the reusable core and why it matters (e.g. every technique
% is scored through the same checker/evaluator, so they are directly comparable).
<CORE ROUTINE> is the reusable core: <how the other stages feed it and why it
makes techniques comparable>.

% =====================================================================
% PER-STAGE / PER-COMPONENT SECTIONS
% Repeat one \section per stage or component listed in the Overview. Each:
%   * prose describing WHAT the stage does and the rules it must respect
%     (tie hard rules back to the requirement IDs [RQxx] where relevant);
%   * a \begin{verbatim} pseudocode block that stands on its own;
%   * a one-line complexity note in the stage's size symbols.
% For an MILP-formulation doc, a "stage" may instead be a constraint FAMILY or
% modeling block; keep the same shape (prose + a precise statement + a note).
% =====================================================================

\section{Stage 1 --- <STAGE 1 NAME>}
% Prose: the stage's job, its inputs/outputs, and every hard rule it enforces
% (cite [RQxx] when a rule comes from a requirement). Say what is FIXED here vs
% what remains a lever for later stages.
<WHAT STAGE 1 DOES and the rules it respects>

\begin{verbatim}
<STAGE 1 PSEUDOCODE>
  # self-contained: inputs, the loop/decision, what it emits.
  # use the algorithm's own names; no source-language specifics required.
\end{verbatim}

% Complexity of this stage alone, in the size symbols you define in \S Complexity.
This runs in $O(<...>)$ (<define the symbols used>).

\section{Stage 2 --- <STAGE 2 NAME>}
<WHAT STAGE 2 DOES>

\begin{verbatim}
<STAGE 2 PSEUDOCODE>
\end{verbatim}

% Duplicate this \section block for each remaining stage/component.

\section{Core routine --- <CORE ROUTINE / OBJECTIVE EVALUATOR NAME>}
% The dedicated section for the routine most other stages call: the objective
% evaluator, cost function, or (for a sequential objective) the discrete-event
% simulator. Give:
%   * its signature and what it returns (the objective value + any schedule /
%     decision it materializes);
%   * its STATE (\paragraph{State.}) if it is a simulator;
%   * its STEP logic (\paragraph{Step.}) as an enumerate, then the pseudocode;
%   * WHY it is correct / checker-valid (\paragraph{...}) — how it satisfies the
%     feasibility rules by construction, so its output passes the P3 checker;
%   * whether it is EXACT or an APPROXIMATION for a fixed input (be explicit;
%     say where optimization actually happens if this routine is greedy).
\texttt{<signature>} is the objective. <one line: what it computes and returns>.

\paragraph{State.} <the mutable state it maintains, if a simulator; else delete>.

\paragraph{Step.} <the unit of work>:
\begin{enumerate}
  \item <selection rule (with tie-break)>.
  \item <action / update>.
  \item <what gets recorded, what advances>.
\end{enumerate}

\begin{verbatim}
<CORE ROUTINE PSEUDOCODE>
\end{verbatim}

\paragraph{Why it is correct / checker-valid.} <how each emitted decision
respects the feasibility rules [RQxx] by construction, so the output passes the
P3 checker>.

\paragraph{Exactness.} <EXACT for its input, or GREEDY/APPROXIMATE — if the
latter, say where the real optimization happens (e.g. the outer search)>.

\section{Complexity}
% Define the size symbols ONCE here (e.g. let N = ..., M = ..., Z = objective
% value / number of steps). Give the cost of one core-routine call, then the
% OVERALL cost as a product across the stages / the outer search. Compare to the
% exact method's worst case where relevant, and note the practical runtime
% regime if measured.
Let <SYMBOL> be <...>, <SYMBOL> be <...>. One <core routine> call costs
$O(<...>)$. The overall cost is
$\approx <...>$ --- <polynomial / exponential; practical runtime note>.

\section{Parameters}
% EVERY tunable knob, with a DEFAULT and its effect. Keep it a simple table.
% If the algorithm has no free parameters, say so in one line and drop the table.
\begin{center}
\begin{tabular}{lll}
\hline
parameter & default & effect \\
\hline
\texttt{<param1>} & <default> & <what it controls> \\
\texttt{<param2>} & <default> & <what it controls> \\
\texttt{<param3>} & <default> & <what it controls> \\
\hline
\end{tabular}
\end{center}

\section{Worked micro-example}
% A tiny instance traced end to end by hand: state the inputs, run each stage /
% a few core-routine steps, and show the resulting objective. Small enough that
% the reader can verify every number. Use a compact table or inline listing.
<INPUTS of the tiny instance>

<TRACE through the stages / steps, ending in the objective value>

\section{Design rationale and limitations}

\paragraph{Rationale.} <WHY this design: which quantity drives the objective,
why the chosen lever/search-space is the right one, and (if applicable) why the
method-review phase recommended this approach>.

\paragraph{Limitations.} <honest list, e.g.>:
(i) <where the routine is greedy/suboptimal and what would tighten it>;
(ii) <what part of the decision space is NOT searched (future work)>;
(iii) <acceptance/stopping-rule or scaling limits and how to relax them>.
